% ---------- Customisation ----------
% Arrow style
\def\arrowcolor{red}
\def\arrowtype{->}          % -stealth, -latex, etc.
\def\arrowthickness{thick}

% Label style
\def\labelfont{\footnotesize}
\def\labelfill{white}
\def\labelprefix{\theta}

% Label offset (in axis units)
\def\labeloffsetx{0}
\def\labeloffsety{0.75}
\def\labeloffsetz{-10}
% -----------------------------------

\begin{figure}[t]
    \centering
    \resizebox{0.85\textwidth}{!}{
        \begin{tikzpicture}
        \begin{axis}[
                % title={$\mathcal{L}(\theta) = w_{1}^2 + w_{2}^2$},
                view={-60}{60},
                xlabel={$w_{1}$},
                ylabel={$w_{2}$},
                zlabel={$\mathcal{L}(w_{1}, w_{2})$},
                colormap/viridis,
                % colorbar,
                % colorbar style={
                %     y dir=reverse,
                %     yticklabel={\pgfmathparse{-\tick}\pgfmathprintnumber{\pgfmathresult}},
                % },
                mesh/ordering=y varies,
                scaled z ticks=false,
        ]

        % ---------- Loss surface ----------
        \addplot3[
            surf,
            shader=flat,
            opacity=0.8,
            faceted color=none,
            point meta=-z,
        ] table {figures/deepl_learning/gradient_descent/simple_grandient_descent/loss_surface.data};

        % ---------- Optimization path ----------
        \pgfplotstableread{figures/deepl_learning/gradient_descent/simple_grandient_descent/optimization_path.data}\datatable
        \pgfplotstablegetrowsof{\datatable}
        \pgfmathtruncatemacro{\numpts}{\pgfplotsretval}
        \pgfmathtruncatemacro{\lastpoint}{\numpts-1}
        \pgfmathtruncatemacro{\lastseg}{\numpts-2}

        % 1. Draw every segment with an arrow
        \foreach \i in {0,...,\lastseg}{
            \pgfmathtruncatemacro{\j}{\i+1}
            \pgfplotstablegetelem{\i}{0}\of\datatable \let\xstart\pgfplotsretval
            \pgfplotstablegetelem{\i}{1}\of\datatable \let\ystart\pgfplotsretval
            \pgfplotstablegetelem{\i}{2}\of\datatable \let\zstart\pgfplotsretval
            \pgfplotstablegetelem{\j}{0}\of\datatable \let\xend\pgfplotsretval
            \pgfplotstablegetelem{\j}{1}\of\datatable \let\yend\pgfplotsretval
            \pgfplotstablegetelem{\j}{2}\of\datatable \let\zend\pgfplotsretval
            \addplot3[
                \arrowcolor,
                \arrowtype,
                \arrowthickness,
                mark=none,
            ] coordinates {
                (\xstart,\ystart,\zstart) (\xend,\yend,\zend)
            };
        }

        % 2. Place a label on every point
        \foreach \k in {0,...,\lastpoint}{
            \pgfmathtruncatemacro{\labelidx}{\k+1}
            \pgfplotstablegetelem{\k}{0}\of\datatable \let\x\pgfplotsretval
            \pgfplotstablegetelem{\k}{1}\of\datatable \let\y\pgfplotsretval
            \pgfplotstablegetelem{\k}{2}\of\datatable \let\z\pgfplotsretval
            % Expand coordinates immediately, but protect font commands
            \edef\templabel{
                \noexpand\node[
                    font=\noexpand\labelfont,
                    fill=\noexpand\labelfill,
                    inner sep=1pt,
                ] at (axis cs: \x+\labeloffsetx,\y+\labeloffsety,\z+\labeloffsetz)
                    {$\labelprefix_{\labelidx}$};
            }
            \templabel
        }

        \end{axis}
        \end{tikzpicture}
    }
    \caption{Gradient descent in a case $\mathcal{L}(w_{1}, w_{2}) = w_{1}^2 + w_{2}^2$.}
    \label{fig:simple_gradient_descent}
\end{figure}